\documentclass[10pt,a4paper]{article}

\usepackage[margin=1.05in]{geometry}
\usepackage{setspace}
\usepackage{amsmath,amssymb,amsthm}
\usepackage{graphicx}
\usepackage{array}
\usepackage{caption}
\usepackage{enumitem}
\usepackage{titlesec}
\usepackage{microtype}
\usepackage{hyperref}
\usepackage{fancyhdr}

\hypersetup{colorlinks=true,linkcolor=black,urlcolor=black,citecolor=black}
\captionsetup{font=small,labelfont=bf,skip=4pt,justification=raggedright}

\titleformat{\section}{\normalsize\bfseries}{}{0em}{}
\titleformat{\subsection}{\small\bfseries}{}{0em}{}
\titleformat{\subsubsection}{\small\itshape}{}{0em}{}
\titlespacing*{\section}{0pt}{1.35em}{0.5em}
\titlespacing*{\subsection}{0pt}{0.95em}{0.35em}

\theoremstyle{plain}
\newtheorem{proposition}{Proposition}
\newtheorem{lemma}{Lemma}
\theoremstyle{definition}
\newtheorem{definition}{Definition}

\setstretch{1.14}
\setlist[itemize]{leftmargin=1.15em,itemsep=0.1em,topsep=0.2em}

\graphicspath{{drawings/}}

\pagestyle{fancy}
\fancyhf{}
\fancyhead[L]{\small\itshape Elements of Position}
\fancyhead[R]{\small\itshape III.\ Meetings}
\fancyfoot[C]{\small\thepage}
\renewcommand{\headrulewidth}{0pt}

\begin{document}

\begin{center}
\vspace*{2.2em}
{\large Elements of Position}\\[1.6em]
{\LARGE\bfseries III.\ Meetings}\\[2.2em]
{\normalsize Justin Erholtz}\\[0.35em]
{\small 12 September 2026}
\end{center}

\vspace{2.4em}

\begin{center}
{\large A's spokes meet B's ratio.}
\end{center}

\vspace{1.8em}

\noindent
Two independent static interiors, overlapped and then walked, are a system.
The system is measured from one reference that does not move: the $1/2$ line.

Part~I built the unit and the $n$-count.
Part~II laid a second process on that count, walked the other way after arrival.
This part asks what happens when those two processes occupy the same drawing.

A freeze is a section, and the rate-$1$ walk makes those sections a system.
The dynamics are not new objects --- they are marks already forced, moving relative to one another while the midline stays.

The first exact marriage is $n=5$: A's own split at five \emph{is} B's ratio.
That identity does not spread to a generic later integer.
Fibonacci docking is a different kind of meeting: approximate, after arrival, shrinking as the $n$-count runs.

\vspace{1.4em}

\noindent
\textit{A note.}
This part cites I and II and stays inside the same questions.
Claims that can be proved here are proved by school geometry.
The drawings are layers on the same architectural set, built by \texttt{Elements of Position - Drawings.py}.

\vspace{0.9em}
\noindent
\textit{A note on names used here.}
Names from I and II stand.
Meeting --- a finished comparison of a mark of A with a mark of B.
Marriage --- an exact meeting, forced by the integers already on the table, not an approximation.
Docking --- a meeting of walks whose error is determined by two arrived integers and shrinks as those integers grow.

\newpage
\tableofcontents
\newpage

% =====================================================================
\section{The common object is the line}
% =====================================================================

The common object is not a circle.
It is the $1/2$ line.

Every local diameter is cut at its midpoint.
Process A's even split lives there.
Process B's pair $\Phi$ and $\Phi^{2}$ lives equally far from it.
The cone of A, if the axis is the $n$-count, has slope $1/2$.
Three facts, one ratio.

Freeze $n$: an $n$-gon, a half-walk, a $\Phi$ mark, an inward copy of radius $(n/2)\Phi$.
All of that is a section.
Start the count and the same marks ride A outward and, after arrival, walk B the other way.

A winds $+2\pi$ per integer; B winds $-2\pi$ per opposite turn after arrival.
Same clock, opposite hands.
At the arrival circle they occupy the same ring and then leave it in opposite angular senses.

They do not have the same length.
If height is walked arc, A from the paper to a finished $N$ climbs $\pi(N^{2}-1)/2$.
B, collapsing at scale $\Phi$ per opposite turn, is bounded by a multiple of $N$.
A grows radius linearly in $n$, so arc piles up quadratically, while B shrinks radius geometrically.
That mismatch is not a drawing error.
It is two processes on one clock.

If the axis is the count $n$, A is the cone $r=n/2$.
B from arrival is $r(n)=(N/2)\Phi^{N-n}$, which is not on that cone.
Same integers, two surfaces.

\begin{proposition}
From a Fibonacci arrival $N=F_k$ already on the table, one opposite turn of B has radius $F_k/(2\varphi)$.
A's freeze at $F_{k-1}$ has radius $F_{k-1}/2$.
The difference is $\tfrac12|F_k\Phi-F_{k-1}|$, which is already determined by those two integers.
\end{proposition}

\begin{proof}
B's first step from $N$ is radius $(N/2)\Phi$.
Substitute $N=F_k$, and compare to the freeze of A at the previous Fibonacci integer already arrived.
\end{proof}

At $N=13$ the miss in radius is about $0.017$; at $N=233$ it is about $0.00096$.
You cannot walk B until you arrive.
When the arrival is a Fibonacci integer already produced by the $n$-count, the first opposite step almost lands on the previous Fibonacci freeze of A.
No extra sequence was imported.

Depth to shrink from diameter $N$ to the paper diameter is $\log_\varphi N$, which is not an integer.
The paper is always a non-integer number of $\Phi$-turns below arrival.
One closes; the cut between $0$ and $1$ does not.

Irreducible interiors do not give B a special radius.
Primality is a fact about A's grammar, and the midline does not care.
The gap from the even cut to the $\Phi$ cut is $n(\Phi-1/2)$ on every diameter: a growing room, same shape.

The system does not surprise itself.
Run $50$ or run $233$ and the same docking error shrinks, the same length mismatch grows, the same $n=5$ identity stays exact, the same $1/2$ line does not move.
That is what rigidity was for: one run is every run.

\begin{figure}[ht]
\centering
\includegraphics[width=0.95\textwidth]{layer_both_01_paper.png}
\caption{Both processes on the paper circle.
A freeze is a section of a live pair.}
\end{figure}

\begin{figure}[ht]
\centering
\includegraphics[width=0.78\textwidth]{layer_both_03_plan.png}
\caption{Plan.
A out with the $n$-count.
B the other way from arrival.}
\end{figure}

\begin{figure}[ht]
\centering
\includegraphics[width=0.95\textwidth]{layer_both_04_obliques.png}
\caption{Obliques.
A up.
B down from arrival, opposite angle, scale $\Phi$.}
\end{figure}

% =====================================================================
\section{The first exact marriage is $n=5$}
% =====================================================================

A's own split at $n=5$ is B's ratio.
The local circle has diameter $5$, radius $R=5/2$, with five equal central angles, each $72^{\circ}$.
That is all A contributes.
B is the pair defined by $1/x=1+x$.

\begin{proposition}
\label{prop:five-quadratic}
Let $\theta=72^{\circ}=2\pi/5$.
Then $\cos72^{\circ}=\Phi/2$, and $x=2\cos72^{\circ}$ satisfies $x^{2}+x-1=0$.
\end{proposition}

\begin{proof}
$2\theta=144^{\circ}=180^{\circ}-\theta$, so $\cos2\theta=-\cos\theta$.
Also $\cos2\theta=2\cos^{2}\theta-1$.
Hence $2c^{2}+c-1=0$ with $c=\cos72^{\circ}$.
The positive root is $(-1+\sqrt5)/4=\Phi/2$.
Set $x=2c$.
Then $x^{2}+x-1=0$, which is $1/x=1+x$.
\end{proof}

One spoke, doubled, produces process B --- no extra map.

The companion is $\cos36^{\circ}=\varphi/2$.
A's side is the chord of one central step; A's diagonal is the chord of two:
\[
\frac{\mathrm{diag}}{\mathrm{side}}=\frac{\sin72^{\circ}}{\sin36^{\circ}}=2\cos36^{\circ}=\varphi.
\]
Every diagonal is cut by another in the ratio $\varphi:\Phi$.
The five intersection points form a smaller regular pentagon of radius $R\Phi^{2}$ --- two inward steps of B.
The first inward copy, radius $R\Phi$, is the step the diagonals skip.

On the diameter, $\Phi$ and $\Phi^{2}$ remain symmetric about the centre.
The vertical diagonal and the $\Phi$-mark are two different $\Phi$-rods from the same origin.
They share the midline and they do not coincide.

Five is the first odd irreducible inversion after $2$.
It is also the first integer whose central angle, doubled, solves $x^{2}+x-1=0$.
The first prime after the even cut, and B's defining equation, arrive together because $360/5=72$ and $2\cdot72$ is the supplement of $72$.

That last sentence is a coincidence of small numbers worth marking, not a claim that every later irreducible arrival will compute B.
At $n=10$ the $10$-gon inherits $36^{\circ}$ as a spoke.
After that A's mesh and B's cut share only centre, diameter, and the $1/2$ line.
The exact computation of B by A's spokes does not repeat at $7$, $8$, $13$, or $16$.

Fibonacci docking was an approximate meeting of walks.
The pentagon is an exact meeting of one freeze.
Both are A against B around $1/2$.

\begin{figure}[ht]
\centering
\includegraphics[width=0.95\textwidth]{layer_n5_meetings.png}
\caption{$n=5$.
Left: spokes, diagonals, $\Phi$ and $\Phi^{2}$ on the diameter, half-walk.
Right: inner pentagon at $R\Phi^{2}$ and first inward copy at $R\Phi$.}
\end{figure}

% =====================================================================
\section{Solved triangles in the five-disk}
% =====================================================================

Every length below is a word in $R=5/2$ and $\Phi$.
The angle alphabet of the disk is $\{18^{\circ},36^{\circ},54^{\circ},72^{\circ},90^{\circ},108^{\circ}\}$.
Two of those are already $\Phi$:
\[
\sin18^{\circ}=\cos72^{\circ}=\frac\Phi2,\qquad
\cos36^{\circ}=\sin54^{\circ}=\frac\varphi2=\frac1{2\Phi}.
\]

\paragraph{Thales.}
Half-walk apex $A=(0,R)$, diameter ends $L=(-R,0)$ and $W=(R,0)$.
Angle at $A$ is $90^{\circ}$ on any $n$.
Legs $R\sqrt2$, hypotenuse $2R=5$.
The half-walk is still only the even cut.
It does not know $\Phi$.

\paragraph{Spoke triangle $54^{\circ}$--$54^{\circ}$--$72^{\circ}$.}
Two radii and one side.
Drop a perpendicular from the centre to the side: two right triangles $36^{\circ}$--$54^{\circ}$--$90^{\circ}$,
\[
\text{apothem}=\frac R{2\Phi},\qquad
s=2R\sin36^{\circ},\qquad
d=s/\Phi.
\]

\paragraph{Similar family $36^{\circ}$--$72^{\circ}$--$72^{\circ}$.}
Two diagonals and a side: $s,d,d$.
A crossing cuts $d$ as $d\Phi=s$ and $d\Phi^{2}=s\Phi$.
Those pieces sum to $d$ because $\Phi+\Phi^{2}=1$.
The same identity that placed $\Phi$ and $\Phi^{2}$ on the diameter now places them on A's diagonal.
The inner side is $s\Phi^{2}$.
The $72^{\circ}$-chord on the first inward copy is $s\Phi$.

\paragraph{Chord of $108^{\circ}$.}
From the left end of the diameter to the first vertex,
\[
|LV_1|=2R\sin54^{\circ}=R/\Phi.
\]

Inventory in one disk: radius $R$; diameter $2R$; half-walk legs $R\sqrt2$; side $s$; diagonal $s/\Phi$; apothem $R/(2\Phi)$; $\Phi$-mark from centre $R(2\Phi-1)$; first inward copy $R\Phi$; inner pentagon $R\Phi^{2}$; short diagonal piece $s\Phi$; long diagonal piece $s$; chord $108^{\circ}$ equal to $R/\Phi$.

Three triangle species share one freeze: the midline holds the even cut, the spokes hold the five-split, and the diagonals hold the similar nest, which is B.

Scale does not matter.
The same inventory at any similar five-disk is the same words times a ratio of parts.
Inside a circle of diameter $10\,000$ there is a similar five-disk whose every length scales by $2000$.
The marriage is the shape, not the size.

% =====================================================================
\section{What high $n$ does and does not do}
% =====================================================================

No limit is taken here.
A large finished integer is only a deeper sample of the same building.

What does not change: rate $1$; midline $1/2$; equal arc $\pi$; $\Phi+\Phi^{2}=1$; opposite helicity; B only after arrival; $n=5$ still computes B and no later generic $n$ joins it.

The two-step chord over the one-step chord at integer $n$ is $2\cos(\pi/n)$.
That equals $\varphi$ only when $\pi/n=36^{\circ}$, hence only at $n=5$.
At $n=10\,000$ the same ratio is already $1.99999961$, which is not $\varphi$.
The marriage does not spread.

A's walked height grows like $N^{2}$, while B's whole opposite walk from that arrival grows like $N$.
Raising the count pins B harder to the top of A's climb.
Depth to the paper diameter remains $\log_\varphi N$, still not an integer.

From a finished $N$, integer depths of B land at diameters $N\Phi^{d}$.
They are not A's integer freezes except at $d=0$.

High $n$ does not dissolve the meetings.
It makes the non-meetings obvious: the exact marriage stays at five, the approximate dockings get closer, and the midline does not move.

\begin{figure}[ht]
\centering
\includegraphics[width=0.95\textwidth]{layer_both_02_interiors.png}
\caption{Local interiors at a freeze.
A's $n$-gon and half-walk against B's first inward copy and the pair $\Phi,\Phi^{2}$.}
\end{figure}

% =====================================================================
\section{What stands}
% =====================================================================

Forced:
the $1/2$ line as the common object of A and B;
opposite helicity on one $n$-count;
length mismatch of the two walks;
docking error $\tfrac12|F_k\Phi-F_{k-1}|$ at Fibonacci arrivals already on the table;
$\cos72^{\circ}=\Phi/2$, so $n=5$ computes B from A's spokes alone;
diag/side $=\varphi$ in the five-disk;
$\Phi+\Phi^{2}=1$ cutting both the diameter and the diagonal;
the angle alphabet $\{18^{\circ},36^{\circ},54^{\circ},72^{\circ},90^{\circ},108^{\circ}\}$ as words in $\Phi$;
$2\cos(\pi/n)=\varphi$ only at $n=5$.

Named:
meeting, marriage, docking.

Not granted:
a second primary ratio;
an exact marriage at a generic later $n$;
a completed Fibonacci sequence;
a walk of B before arrival;
a new object beyond the two processes already built.

The first sentence is still the sentence.
Where is $1$?
At $n=5$ one answer is: on the midline, with A's spokes naming B's cut, and B's cut sitting equally far from that line as its own square.
That is a meeting.
It is not a new unit.

\end{document}
